\chapter*{Abstract}
\addcontentsline{toc}{chapter}{Abstract}

Persuasive dialogue is central to charitable fundraising, yet the
Persuasion-For-Good corpus \citep{wang2019persuasion} -- the largest public
resource of real donation-solicitation dialogues -- ships with only a
binary donation outcome and a small demonstrative sample of strategy labels.
It is therefore unknown, at scale, \emph{which} persuasion strategies a
persuader actually used turn by turn, and consequently unknown which
strategies are associated with a donation once confounds like dialogue
length are controlled for. This project closes that gap end to end: a
41-strategy / 11-category multi-label taxonomy was designed and applied to
all 10{,}600 persuader turns of the corpus (1{,}017 dialogues) under an
LLM-assisted, human-checkpointed annotation process; a multi-label
classifier was built to recognize these strategies automatically; and the
recovered strategy labels were used to model donation outcome.

The classification system evolved through seventeen measured iterations
(v1--v17). Each iteration is a controlled experiment -- a documented
hypothesis, a change scoped to as few variables as possible, deterministic
\emph{sentinel} ensemble members whose scores must reproduce bit-for-bit
between an iteration and its predecessor before any comparison is trusted,
and a decision recorded whether the change survived contact with held-out
data. Several plausible ideas were rejected this way (target-span attention
pooling, prior-initialized classification-head biases, a naive
class-imbalance loss reweighting that improved micro- but not macro-F1);
two were kept (an architecture-dependent context-width effect, and an
empirical-Bayes threshold-shrinkage decision rule that this project
contributes). The most consequential single change was reframing the task
itself: rather than predicting all 41 fine-grained strategies and reading
off the 11 coarse categories as a byproduct, the final system (v15--v17)
makes the 11-category task the primary training and decoding objective,
with the 41-strategy task retained as an auxiliary head whose gradient
demonstrably regularizes the rare categories. On the held-out test split,
this raises category-level micro-F1 to \textbf{0.706} and macro-F1 to
\textbf{0.564} -- within \SI{0.5}{pp} of the corpus's own cross-annotator
agreement ceiling on the same task ($\kappa\text{-analog}$ consensus bound
0.793), meaning the remaining error is close to the noise floor of the
human labels themselves, concentrated almost entirely in three
categories that make up 3.4\% of all positive labels.

Feeding these recovered strategy labels into a McFadden pseudo-$R^2$
logistic-regression outcome model raises cross-validated explained variance
from 0.12 (categories alone) to \textbf{0.165} (with strategy co-occurrence
features) on the same 1{,}017 dialogues -- roughly double the 0.015--0.08
range reported by prior single-label work on a related corpus, though this
project is explicit that the comparison is confounded by more than the
labeling scheme and should not be over-read. A companion donation-\emph{intent}
classifier built by the same team, using a dual-stream cross-attention
architecture over persuader/persuadee turns, separately reaches
\textbf{89.5\%} accuracy and \textbf{0.862} macro-F1 on binary donation
intent, and the same architecture transfers to a structurally similar
buyer--seller negotiation domain at \textbf{97.3\%} accuracy.

Beyond the numbers, this report documents a reproducible experimental
methodology for iterating on a multi-label classifier under a hard compute
budget (a single free Kaggle GPU session per run): GPU-free smoke tests that
exercise every code path on a data subsample before committing to a
multi-hour run, redundant training across two Kaggle accounts to survive
weekly GPU-quota exhaustion, and a statistically honest reporting
convention -- cluster bootstrap confidence intervals over dialogues rather
than turns, dual micro-/macro-optimal decision thresholds fitted by
grouped cross-validation, and an explicit annotation-reliability ceiling
reported alongside every headline score so that a result is never presented
as better, or worse, than the data can actually support.
